\begin{abstract}
Perplexity and downstream performance can diverge after compression, but
post-pruning recovery is still commonly summarized by a loss curve and a small
set of endpoints. We present, to our knowledge, the first \emph{multi-interface
recovery audit} that jointly follows same-source perplexity, two MMLU scoring
interfaces, and zero-shot closed-book QA while pairing them with an intact
continued-pretraining control, a random null floor, and an alternative
same-depth construction. On OLMo-2-7B, a 16-layer keep14+fresh2 model improves
through 200k steps to PPL 10.561, yet reaches only .319 MMLU versus .605 for the
intact base and remains 11.6--34.2 points behind on PopQA, TriviaQA, and
NQ-open. Complete-option scoring raises keep14 by 6.48 points, but raises a
fully random 16-layer model by 11.28 points to the same .360 content-score
floor, showing why an interface gain needs a null baseline. At the same nominal
depth and 200k-step budget, ShortGPT reaches PPL 9.780 and MMLU .474, revealing
a 15.5-point construction gap over keep14. The intact full32 branch remains much
closer to the base over its available 25k horizon. Together, these findings
establish a practical measurement lesson: recovery should be audited across
likelihood, target behavior, evaluation interface, construction, budget, and
run-level uncertainty rather than certified by perplexity improvement alone.
\end{abstract}
